\chapter{Personel, ślady cyfrowe, metadane dokumentów i repozytoria}\label{ch:personel-slady}

% ==============================================================================
% ROZDZIAŁ 5 - WSKAZÓWKI I TO-DO DLA AGENTA PISZĄCEGO:
% Rola: Osoba 2 (Personel, adresacja e-mail, ExifTool, SOCMINT, GitLab/GitHub, bazy wycieków).
% Cel: Profilowanie stacji roboczych pod canary token, identyfikacja personelu, telemetria naruszeń.
% Anonimizacja: Rygorystyczne stosowanie \anon{Rola_Jednostka_Inicjały} we wszystkich tabelach i tekście.
% Objętość docelowa: ~7-8 stron.
% ==============================================================================

\section{Struktura adresacji e-mail i punkty kontaktowe}\label{sec:struktura-email}

\subsection{Schematy adresacji pocztowej pracowników i studentów}\label{subsec:schematy-email}
% TO-DO:
% 1. Przedstawić schemat poczty pracowników i doktorantów: <imie>.<nazwisko>@inf.ethz.ch oraz techniczny alias nethz.
% 2. Przedstawić schemat poczty studentów: <nethz_login>@student.ethz.ch (8-znakowy identyfikator).
% 3. Wyjaśnić znaczenie wywiadowcze schematów dla generowania słowników pod kątem spear phishingu.

\subsection{Katalogi osobowe i skrzynki funkcyjne}\label{subsec:ksiazki-adresowe}
% TO-DO:
% 1. Opisać ogólnouczelnianą wyszukiwarkę ETH People Directory oraz katalog wydziałowy inf.ethz.ch/people.html.
% 2. Zidentyfikować skrzynki funkcyjne o wysokim zaufaniu: helpdesk@inf.ethz.ch, studies@inf.ethz.ch, zisc-info@inf.ethz.ch.
% 3. Wskazać ryzyko podszywania się pod adresy funkcyjne w kampaniach socjotechnicznych.

\section{Ekstrakcja i analiza metadanych dokumentów publicznych}\label{sec:metadane-plikow}

\subsection{Dobór próby dokumentów i metodyka ekstrakcji}\label{subsec:pozyskiwanie-dokumentow}
% TO-DO:
% 1. Pasywne pobranie reprezentatywnej próby dokumentów (formularze ISG, regulaminy laboratoriów, wnioski grantowe).
% 2. Zastosowanie narzędzia ExifTool do lokalnej ekstrakcji metadanych bez ingerencji w serwery uczelni.

\subsection{Analiza metadanych środowiskowych (ExifTool)}\label{subsec:analiza-exiftool}
% TO-DO:
% 1. Tabela zanonimizowanych metadanych dla próby plików (z identyfikatorami \anon{Inzynier_ISG_AB},
%    \anon{Doktorant_Sys_MK}, \anon{Sekretariat_CAB_JS}, \anon{Kierownik_Kat_GH} itd.):
%    - Oprogramowanie tworzące (np. TeX Live, Quartz PDFContext, MS Word for Mac, Acrobat Distiller, LibreOffice).
%    - Ujawnione ścieżki lokalne i sieciowe (np. katalogi domowe użytkowników, zasoby współdzielone SMB/NFS).
% 2. Wnioski: heterogeniczność stacji roboczych (Linux/macOS/Windows) wykorzystana w Scenariuszu 2 (Rozdział 6)
%    do przygotowania nieszkodliwego nośnika PDF z canary tokenem.

\section{Profilowanie w mediach społecznościowych (SOCMINT)}\label{sec:socmint-tozsamosci}

\subsection{Oficjalne kanały komunikacji wydziałowej}\label{subsec:oficjalne-social-media}
% TO-DO:
% 1. Przeanalizować kanały instytucjonalne: profil LinkedIn (>20 tys. obserwujących), X (@ETH_DINFK), YouTube.
% 2. Wskazać pozyskane informacje: relacje z konferencji, zdjęcia wnętrz CAB, zapowiedzi grantów i harmonogramy nieobecności.

\subsection{Profile zawodowe i naukowe personelu}\label{subsec:profile-zawodowe}
% TO-DO:
% 1. Zanalizować profile personelu na LinkedIn, ResearchGate, Google Scholar i dblp (zanonimizowane przez \anon{...}).
% 2. Skategoryzować profile: kadra badawcza (np. PyTorch, LLVM, C++, Rust, FPGA) vs kadra inżynieryjna (np. Kubernetes, Docker, Slurm, Ceph).
% 3. Wykorzystać profilowanie techniczne do personalizacji treści spear phishingowych.

\subsection{Korelacja pseudonimów i kont (WhatsMyName)}\label{subsec:korelacja-whatsmyname}
% TO-DO:
% 1. Zastosować narzędzie whatsmyname.app do pasywnej weryfikacji unikalnych pseudonimów personelu D-INFK.
% 2. Zidentyfikować powiązane konta na platformach deweloperskich i forach (np. Stack Overflow, GitHub, Reddit, Keybase).
% 3. Zidentyfikować wycieki operacyjne (OPSEC leakage) w pytaniach zadawanych na forach technicznych.

\section{Aktywność w otwartych repozytoriach kodu źródłowego}\label{sec:repozytoria-kodu}

\subsection{Publiczne repozytoria GitLab D-INFK oraz GitHub}\label{subsec:projekty-gitlab-github}
% TO-DO:
% 1. Przeanalizować publiczną instancję gitlab.inf.ethz.ch (kursy, szablony laboratoriów studenckich).
% 2. Przeanalizować organizacje badawcze D-INFK na platformie GitHub (np. eth-sri, spcl, systems-ethz).

\subsection{Analiza historii commitów i potoków CI/CD}\label{subsec:analiza-commitow}
% TO-DO:
% 1. Zbadać metadane commitów: korelacja służbowych adresów @inf.ethz.ch z prywatnymi kontami deweloperów.
% 2. Przeanalizować pliki CI/CD (.gitlab-ci.yml, GitHub Actions): obrazy Docker, tagi runnerów wewnętrznych klastrów.
% 3. Przeprowadzić pasywny audyt secret scanning: ocenić higienę kodu (potwierdzenie High Security Posture przy braku poświadczeń).

\section{Pasywna weryfikacja ekspozycji w wyciekach danych}\label{sec:wycieki-danych}

\subsection{Analiza ekspozycji domenowej w bazach naruszeń}\label{subsec:bazy-naruszen}
% TO-DO:
% 1. Wykorzystać darmową i legalną telemetrię naruszeń: Intelligence X (intelx.io), raporty infostealerów (np. Hudson Rock), LeakIX.
% 2. Zestawić statystyki: dominacja starych incydentów zewnętrznych (serwisy konferencyjne i edukacyjne).
% 3. Zinterpretować skalę ekspozycji domeny inf.ethz.ch pod kątem dojrzałości higieny haseł (High Security Posture).

\subsection{Ocena ryzyka ataków na poświadczenia}\label{subsec:ryzyko-poswiadczen}
% TO-DO:
% 1. Przeprowadzić modelowanie ryzyka: wykazać nieskuteczność credential stuffingu i password sprayingu.
% 2. Zastosować zasadę High Security Posture: egzekwowanie MFA na bramie Switch edu-ID, rate limiting, blokady kont.
% 3. Wnioski dla Rozdziału 6: konieczność zastosowania zaawansowanego ataku AitM (Evilginx2) zamiast prostych ataków słownikowych.
